\documentclass[11pt]{article}
\usepackage[utf8]{inputenc}
\usepackage[T1]{fontenc}
\usepackage[italian]{babel}
\usepackage{geometry}
\geometry{a4paper, margin=1in}
\usepackage{xcolor}
\usepackage{listings}
\usepackage{hyperref}

\usepackage{tcolorbox}
\tcbuselibrary{breakable}
\begin{document}

\begin{tcolorbox}[colback=blue!5!white,colframe=blue!75!black,title=\textbf{DOMANDA}, breakable]
\begin{itemize}
    \item Aprire e studiare con Wireshark il file PCAP contenuto nello ZIP
    \begin{itemize}
        \item Quale protocollo di livello trasporto utilizza MQTT? Quali sono le porte?
        \item Trovare le fasi di publish e subscribe. Quante connessioni \texttt{TCP} apre un subscriber? Quante connessioni \texttt{TCP} apre un publisher?
    \end{itemize}
\end{itemize}
\end{tcolorbox}

\begin{tcolorbox}[colback=green!5!white,colframe=green!75!black,title=\textbf{RISPOSTA}, breakable]
\textbf{Soluzione esercizio MQTT - es42 (Analisi PCAP)}

\textbf{Domanda 1:} Quale protocollo di livello trasporto utilizza MQTT? Quali sono le porte?\\
\textbf{Risposta 1:}\\
MQTT utilizza il protocollo di livello trasporto \texttt{TCP}.\\
La porta standard utilizzata (come si vede dal file PCAP) è la 1883 per le connessioni in chiaro (non cifrate). Qualora si utilizzasse MQTT over TLS (comunicazione cifrata), la porta standard sarebbe la 8883.

\textbf{Domanda 2:} Trovare le fasi di publish e subscribe. Quante connessioni \texttt{TCP} apre un subscriber? Quante connessioni \texttt{TCP} apre un publisher?\\
\textbf{Risposta 2:}\\
Entrambi i client (publisher e subscriber) instaurano esattamente 1 sola connessione \texttt{TCP} con il broker per svolgere il proprio lavoro.
\begin{itemize}
    \item \textbf{Il Subscriber apre 1 singola connessione \texttt{TCP} persistente:} esegue il 3-way handshake, invia il messaggio MQTT di "Connect Command", riceve il "Connect Ack", dopodiché invia la "Subscribe Request" e mantiene aperta quella stessa connessione \texttt{TCP} per rimanere in ascolto dei futuri messaggi.
    \item \textbf{Il Publisher apre anch'esso 1 singola connessione \texttt{TCP}:} esegue l'handshake, si connette ("Connect Command"), invia il dato ("Publish Message") e al termine invia una richiesta di disconnessione ("Disconnect Req"), chiudendo subito dopo la singola connessione \texttt{TCP} utilizzata.
\end{itemize}

\textbf{Domanda 3:} Documentare i comandi utilizzati per aprire e analizzare sia da interfaccia grafica che utilizzando la CLI.\\
\textbf{Risposta 3:}

\textbf{--- Analisi da Interfaccia Grafica (Wireshark) ---}\\
Puoi aprire il file lanciando da terminale:
\begin{lstlisting}[language=bash, basicstyle=\ttfamily\small]
wireshark mqtt-localhost.pcapng
\end{lstlisting}
Una volta aperta l'interfaccia di Wireshark, per isolare immediatamente e studiare le fasi esatte di publish e subscribe ignorando il rumore di fondo, è sufficiente scrivere nella barra dei filtri in alto (Display Filter) la parola chiave:
\texttt{mqtt}
(e poi premere Invio).

\textbf{--- Analisi da riga di comando (CLI) ---}\\
Se vuoi analizzare il PCAP interamente da riga di comando usando TShark (la versione CLI di Wireshark), il comando per filtrare solo i pacchetti MQTT è:
\begin{lstlisting}[language=bash, basicstyle=\ttfamily\small]
tshark -r mqtt-localhost.pcapng -Y "mqtt"
\end{lstlisting}
In alternativa, se tshark non è disponibile, puoi usare il classico \texttt{tcpdump} per visualizzare tutti i pacchetti \texttt{TCP} scambiati sulla porta 1883:
\begin{lstlisting}[language=bash, basicstyle=\ttfamily\small]
tcpdump -r mqtt-localhost.pcapng -n port 1883
\end{lstlisting}
\end{tcolorbox}

\end{document}
